%% Basic Syntax and Semantics
\chapter{Basic Syntax and Semantics}
\label{ch:syntax}

This section defines the core syntactic forms of Lambdust, extending R7RS-small with new constructs while maintaining complete backward compatibility. The syntax is organized into expressions, which are the primary syntactic category.

\section{Syntactic Categories}
\label{sec:syntactic-categories}

Lambdust programs consist of expressions organized into the following syntactic categories:

\begin{grammar}
\production{\var{program}}{\produces}{\arbno{\var{import-declaration}} \arbno{\var{expression}}}{}
\\
\production{\var{expression}}{\produces}{\var{constant}}{}
\production{}{\alternative}{\var{variable}}{}
\production{}{\alternative}{\var{procedure-call}}{}
\production{}{\alternative}{\var{lambda-expression}}{}
\production{}{\alternative}{\var{conditional}}{}
\production{}{\alternative}{\var{assignment}}{}
\production{}{\alternative}{\var{derived-expression}}{}
\production{}{\alternative}{\var{macro-use}}{}
\production{}{\alternative}{\var{type-expression}}{}
\production{}{\alternative}{\var{effect-expression}}{}
\production{}{\alternative}{\var{concurrency-expression}}{}
\production{}{\alternative}{\var{foreign-expression}}{}
\end{grammar}

\section{Constants and Literals}
\label{sec:constants}

\indexentry{constant}
\indexentry{literal}

Constants are self-evaluating expressions:

\begin{grammar}
\production{\var{constant}}{\produces}{\var{boolean} \alternative{} \var{number}}{}
\production{}{\alternative}{\var{character} \alternative{} \var{string}}{}
\production{}{\alternative}{\var{keyword}}{}
\end{grammar}

\begin{example}[Constants]
\begin{lambdustcode}
#t              ; Boolean true
42              ; Integer  
3.14            ; Real number
1/3             ; Rational number
3+4i            ; Complex number
#\lambda        ; Character
"Hello"         ; String
#:type          ; Keyword
\end{lambdustcode}
\end{example}

\section{Variables}
\label{sec:variables}

\indexentry{variable}

Variables are identifiers that refer to bound values:

\begin{grammar}
\production{\var{variable}}{\produces}{\var{identifier}}{}
\end{grammar}

Variable references are resolved lexically according to R7RS scoping rules.

\section{Procedure Calls}
\label{sec:procedure-calls}

\indexentry{procedure call}
\indexentry{application}

Procedure calls apply a procedure to arguments:

\begin{grammar}
\production{\var{procedure-call}}{\produces}{(\var{operator} \arbno{\var{operand}})}{}
\production{\var{operator}}{\produces}{\var{expression}}{}
\production{\var{operand}}{\produces}{\var{expression}}{}
\end{grammar}

\begin{example}[Procedure Calls]
\begin{lambdustcode}
(+ 1 2 3)                   ; Add three numbers
(cons 'a '(b c))           ; Cons onto list
((lambda (x) (+ x 1)) 5)   ; Apply lambda
(f x y z)                  ; Call f with three arguments
\end{lambdustcode}
\end{example}

\section{Lambda Expressions}
\label{sec:lambda}

\indexentry{lambda}
\indexentry{procedure}

Lambda expressions create procedures:

\begin{grammar}
\production{\var{lambda-expression}}{\produces}{(lambda \var{formals} \arbno{\var{expression}})}{}
\\
\production{\var{formals}}{\produces}{(\arbno{\var{identifier}})}{}
\production{}{\alternative}{(\arbno{\var{identifier}} . \var{identifier})}{}
\production{}{\alternative}{\var{identifier}}{}
\end{grammar}

\begin{example}[Lambda Expressions]
\begin{lambdustcode}
(lambda (x) (+ x 1))           ; Single parameter
(lambda (x y) (* x y))         ; Multiple parameters  
(lambda args (length args))    ; Variable arguments
(lambda (x . rest) 
  (cons x rest))               ; Fixed + variable arguments
\end{lambdustcode}
\end{example}

\section{Conditionals}
\label{sec:conditionals}

\indexentry{if}
\indexentry{conditional}

Conditional expressions perform branches based on test values:

\begin{grammar}
\production{\var{conditional}}{\produces}{(if \var{test} \var{consequent} \var{alternate})}{}
\production{}{\alternative}{(if \var{test} \var{consequent})}{}
\end{grammar}

\begin{example}[Conditionals]
\begin{lambdustcode}
(if (> x 0) 'positive 'non-positive)    ; Full conditional
(if (null? lst) 'empty)                 ; No alternate
\end{lambdustcode}
\end{example}

\section{Assignment}
\label{sec:assignment}

\indexentry{set!}
\indexentry{assignment}

Assignment expressions modify variable bindings:

\begin{grammar}
\production{\var{assignment}}{\produces}{(set! \var{variable} \var{expression})}{}
\end{grammar}

\begin{example}[Assignment]
\begin{lambdustcode}
(set! x 42)        ; Assign 42 to x
(set! counter (+ counter 1))  ; Increment counter
\end{lambdustcode}
\end{example}

\section{Derived Expressions}
\label{sec:derived}

\indexentry{derived expression}

Derived expressions are syntactic sugar that can be expressed in terms of more primitive forms.

\section{Conditionals}
\label{sec:derived-conditionals}

\begin{syntaxspec}{(cond \arbno{\var{clause}})}
Multi-way conditional with optional else clause.
\begin{lambdustcode}
(cond ((< x 0) 'negative)
      ((= x 0) 'zero)  
      (else 'positive))
\end{lambdustcode}
\end{syntaxspec}

\begin{syntaxspec}{(case \var{key} \arbno{\var{clause}})}
Dispatch on datum equality.
\begin{lambdustcode}
(case op
  ((+) (+ x y))
  ((-) (- x y))
  (else (error "unknown op")))
\end{lambdustcode}
\end{syntaxspec}

\section{Boolean Operations}
\label{sec:boolean-ops}

\begin{syntaxspec}{(and \arbno{\var{expression}})}
Short-circuit logical and.
\begin{lambdustcode}
(and (> x 0) (< x 100))    ; Check range
\end{lambdustcode}
\end{syntaxspec}

\begin{syntaxspec}{(or \arbno{\var{expression}})}
Short-circuit logical or.
\begin{lambdustcode}
(or (null? lst) (pair? lst))   ; Check list type
\end{lambdustcode}
\end{syntaxspec}

\section{Binding Forms}
\label{sec:binding}

\begin{syntaxspec}{(let (\arbno{\var{binding}}) \arbno{\var{expression}})}
Sequential local binding.
\begin{lambdustcode}
(let ((x 1) (y 2))
  (+ x y))
\end{lambdustcode}
\end{syntaxspec}

\begin{syntaxspec}{(let* (\arbno{\var{binding}}) \arbno{\var{expression}})}
Sequential binding with each binding visible to later ones.
\begin{lambdustcode}
(let* ((x 1) (y (+ x 1)))
  (+ x y))
\end{lambdustcode}
\end{syntaxspec}

\begin{syntaxspec}{(letrec (\arbno{\var{binding}}) \arbno{\var{expression}})}
Recursive binding for mutual recursion.
\begin{lambdustcode}
(letrec ((even? (lambda (n) 
                  (if (= n 0) #t (odd? (- n 1)))))
         (odd? (lambda (n)
                 (if (= n 0) #f (even? (- n 1))))))
  (even? 42))
\end{lambdustcode}
\end{syntaxspec}

\section{Sequencing}
\label{sec:sequencing}

\begin{syntaxspec}{(begin \arbno{\var{expression}})}
Sequential evaluation returning last result.
\begin{lambdustcode}
(begin
  (display "Hello ")
  (display "World")
  (newline))
\end{lambdustcode}
\end{syntaxspec}

\section{Iteration}
\label{sec:iteration}

\begin{syntaxspec}{(do (\arbno{\var{variable}} \var{init} \var{step}) (\var{test} \arbno{\var{result}}) \arbno{\var{command}})}
General iteration construct.
\begin{lambdustcode}
(do ((i 0 (+ i 1))
     (sum 0 (+ sum i)))
    ((>= i 10) sum))    ; Sum 0 to 9
\end{lambdustcode}
\end{syntaxspec}

\section{Definitions}
\label{sec:definitions}

\indexentry{define}
\indexentry{definition}

Definitions bind identifiers to values or procedures:

\begin{grammar}
\production{\var{definition}}{\produces}{(define \var{variable} \var{expression})}{}
\production{}{\alternative}{(define (\var{variable} \arbno{\var{identifier}}) \arbno{\var{expression}})}{}
\production{}{\alternative}{(define (\var{variable} . \var{identifier}) \arbno{\var{expression}})}{}
\end{grammar}

\begin{example}[Definitions]
\begin{lambdustcode}
(define pi 3.14159)                    ; Variable definition

(define (square x) (* x x))            ; Procedure definition

(define (list-of-squares . numbers)    ; Variadic procedure
  (map square numbers))
\end{lambdustcode}
\end{example}

\section{Quoting and Quasiquoting}
\label{sec:quoting}

\indexentry{quote}
\indexentry{quasiquote}

\section{Quote}
\label{sec:quote}

Quote prevents evaluation and returns literal data:

\begin{grammar}
\production{\var{quotation}}{\produces}{(quote \var{datum})}{}
\production{}{\alternative}{\textquotesingle{}\var{datum}}{}
\end{grammar}

\begin{example}[Quotation]
\begin{lambdustcode}
(quote (a b c))     ; List of symbols
'(a b c)            ; Abbreviated form
'symbol             ; Symbol  
'42                 ; Number (self-evaluating)
\end{lambdustcode}
\end{example}

\section{Quasiquote}
\label{sec:quasiquote}

Quasiquote allows selective evaluation within quoted structure:

\begin{grammar}
\production{\var{quasiquotation}}{\produces}{(quasiquote \var{template})}{}
\production{}{\alternative}{\textasciigrave{}\var{template}}{}
\production{\var{template}}{\produces}{\var{datum}}{}
\production{}{\alternative}{(unquote \var{expression})}{}
\production{}{\alternative}{,\var{expression}}{}
\production{}{\alternative}{(unquote-splicing \var{expression})}{}  
\production{}{\alternative}{,@\var{expression}}{}
\end{grammar}

\begin{example}[Quasiquotation]
\begin{lambdustcode}
`(list ,(+ 1 2) 4)          ; Evaluates to (list 3 4)
`(a ,@'(b c) d)             ; Evaluates to (a b c d)
`((f x) = ,(f x))           ; Template with evaluation
\end{lambdustcode}
\end{example}

\section{Type Expressions}
\label{sec:type-expressions}

\indexentry{type expression}

Lambdust extends syntax with type expressions for gradual typing:

\begin{grammar}
\production{\var{type-expression}}{\produces}{(:: \var{expression} \var{type})}{}
\production{}{\alternative}{(\var{expression} :: \var{type})}{}
\\
\production{\var{type}}{\produces}{\var{base-type}}{}
\production{}{\alternative}{\var{compound-type}}{}
\production{}{\alternative}{\var{dependent-type}}{}
\\
\production{\var{base-type}}{\produces}{Dynamic \alternative{} Number \alternative{} String}{}
\production{}{\alternative}{Boolean \alternative{} Symbol \alternative{} Character}{}
\\
\production{\var{compound-type}}{\produces}{(\code{->} \arbno{\var{type}} \var{type})}{}
\production{}{\alternative}{(List \var{type})}{}
\production{}{\alternative}{(Vector \var{type})}{}
\production{}{\alternative}{(\var{type} \alternative{} \var{type})}{}
\end{grammar}

\begin{example}[Type Expressions]
\begin{lambdustcode}
(:: 42 Number)                    ; Type ascription
(:: (+ 1 2) Natural)             ; Refined type
(:: factorial (-> Natural Natural))  ; Function type
(:: lst (List Number))           ; Parameterized type
\end{lambdustcode}
\end{example}

\section{Effect Expressions}  
\label{sec:effect-expressions}

\indexentry{effect}
\indexentry{perform}
\indexentry{handle}

Effect expressions provide structured side effect management:

\begin{grammar}
\production{\var{effect-expression}}{\produces}{(perform \var{effect-name} \arbno{\var{expression}})}{}
\production{}{\alternative}{(handle \var{expression} \var{handler})}{}
\production{}{\alternative}{(with-effect \var{effect-name} \var{expression})}{}
\\
\production{\var{handler}}{\produces}{(\var{effect-name} \arbno{\var{operation-clause}})}{}
\production{\var{operation-clause}}{\produces}{(\var{operation} \var{parameters} \var{body})}{}
\end{grammar}

\begin{example}[Effect Expressions]
\begin{lambdustcode}
;; Perform an effect
(perform State get)
(perform State set 42)

;; Handle effects
(handle 
  (begin 
    (perform State set 10)
    (perform State get))
  (State 
    ((get k) (k current-state))
    ((set new-state k) 
     (set! current-state new-state)
     (k 'unit))))
\end{lambdustcode}
\end{example}

\section{Concurrency Expressions}
\label{sec:concurrency-expressions}

\indexentry{spawn}
\indexentry{send}
\indexentry{receive}

Concurrency expressions support actor-based parallel programming:

\begin{grammar}
\production{\var{concurrency-expression}}{\produces}{(spawn \var{expression})}{}
\production{}{\alternative}{(send \var{expression} \var{expression})}{}
\production{}{\alternative}{(receive \arbno{\var{message-clause}})}{}
\production{}{\alternative}{(future \var{expression})}{}
\production{}{\alternative}{(await \var{expression})}{}
\\
\production{\var{message-clause}}{\produces}{(\var{pattern} \var{expression})}{}
\end{grammar}

\begin{example}[Concurrency Expressions]
\begin{lambdustcode}
;; Spawn an actor
(define pid (spawn (lambda () (server-loop))))

;; Send a message
(send pid '(request 42))

;; Receive messages with pattern matching
(receive
  (('request data) 
   (handle-request data))
  (('stop)
   (shutdown-server))
  ((other) 
   (log-unknown-message other)))

;; Futures for async computation
(let ((f (future (expensive-computation))))
  (do-other-work)
  (await f))    ; Wait for result
\end{lambdustcode}
\end{example}

\section{Foreign Function Interface}
\label{sec:foreign-syntax}

\indexentry{foreign-call}
\indexentry{foreign-data}

Foreign function interface expressions enable C library integration:

\begin{grammar}
\production{\var{foreign-expression}}{\produces}{(foreign-call \var{name} \var{type} \arbno{\var{expression}})}{}
\production{}{\alternative}{(foreign-data \var{name} \var{type})}{}
\production{}{\alternative}{(with-capability \var{capability} \var{expression})}{}
\end{grammar}

\begin{example}[Foreign Function Interface]
\begin{lambdustcode}
;; Declare foreign function  
(foreign-declare "math.h"
  (sqrt :: (-> Float64 Float64)))

;; Call foreign function with capability
(with-capability 'math-library
  (lambda ()
    (foreign-call sqrt (-> Float64 Float64) 16.0)))

;; Access foreign data
(foreign-data errno (-> Unit Int32))
\end{lambdustcode}
\end{example}

\section{Syntax Extensions}
\label{sec:syntax-extensions}

\indexentry{syntax-rules}
\indexentry{syntax-case}

Lambdust supports R7RS hygenic macros plus extensions:

\begin{grammar}
\production{\var{macro-definition}}{\produces}{(define-syntax \var{name} \var{transformer})}{}
\\
\production{\var{transformer}}{\produces}{(syntax-rules \var{literals} \arbno{\var{rule}})}{}
\production{}{\alternative}{(syntax-case \var{expression} \var{literals} \arbno{\var{clause}})}{}
\production{}{\alternative}{(lambda (\var{stx}) \var{body})}{}
\end{grammar}

\begin{example}[Syntax Extensions]
\begin{lambdustcode}
;; Simple macro with syntax-rules
(define-syntax when
  (syntax-rules ()
    ((when test body ...)
     (if test (begin body ...)))))

;; Pattern matching macro
(define-syntax match
  (syntax-rules ()
    ((match expr 
       (pattern body ...) ...)
     (match-implementation expr 
       (list (cons pattern (lambda () body ...)) ...)))))

;; Usage
(when (> x 0) 
  (display "positive")
  (newline))

(match lst
  (() 'empty)
  ((x) 'singleton)
  ((x y) 'pair)  
  (other 'longer))
\end{lambdustcode}
\end{example}

\section{Evaluation Order}
\label{sec:evaluation}

\indexentry{evaluation order}

Lambdust follows R7RS evaluation order:

\begin{itemize}
\item Procedure calls evaluate operator first, then operands left to right
\item Conditional test is evaluated before branches
\item Assignment evaluates expression before storing
\item Let bindings evaluate initializers left to right
\item Lambda expressions are not evaluated (they create closures)
\end{itemize}

\section{Tail Calls}
\label{sec:tail-calls}

\indexentry{tail call}
\indexentry{proper tail recursion}

Lambdust provides proper tail recursion as required by R7RS:

\begin{definition}[Tail Position]
An expression is in tail position if its value is returned directly as the value of the enclosing procedure without further computation.
\end{definition}

\begin{example}[Tail Recursion]
\begin{lambdustcode}
;; Tail recursive factorial
(define (factorial n acc)
  (if (<= n 1) 
      acc                           ; Tail position
      (factorial (- n 1) (* n acc)))) ; Tail call

;; Not tail recursive  
(define (factorial-bad n)
  (if (<= n 1)
      1
      (* n (factorial-bad (- n 1))))) ; Not in tail position
\end{lambdustcode}
\end{example}

\section{Continuations}
\label{sec:continuations}

\indexentry{call/cc}
\indexentry{continuation}

Lambdust supports full continuations via \code{call-with-current-continuation}:

\begin{procspec}{call-with-current-continuation}
\textbf{call-with-current-continuation} \var{proc} \goesto{} \var{obj}

Captures the current continuation and passes it as a procedure to \var{proc}. The continuation procedure, when called, abandons the current computation and returns its argument as the result of the \code{call-with-current-continuation} call.
\end{procspec}

\begin{example}[Continuations]
\begin{lambdustcode}
;; Early return from loop
(define (find-positive lst)
  (call/cc 
    (lambda (return)
      (for-each 
        (lambda (x) 
          (when (> x 0) (return x)))
        lst)
      #f)))

;; Non-local exit
(define (deep-computation x)
  (call/cc
    (lambda (escape)
      (let ((result (complex-calculation x)))
        (if (valid? result)
            result
            (escape 'error))))))
\end{lambdustcode}
\end{example}

\section{Multiple Values}
\label{sec:multiple-values}

\indexentry{values}
\indexentry{call-with-values}

Lambdust supports multiple value returns:

\begin{procspec}{values \arbno{\var{obj}}}
Returns multiple values.
\end{procspec}

\begin{procspec}{call-with-values \var{producer} \var{consumer}}
Calls \var{producer} with no arguments and passes all returned values to \var{consumer}.
\end{procspec}

\begin{example}[Multiple Values]
\begin{lambdustcode}
;; Return multiple values
(define (quotient-remainder x y)
  (values (quotient x y) (remainder x y)))

;; Consume multiple values  
(call-with-values 
  (lambda () (quotient-remainder 17 5))
  (lambda (q r) 
    (display "Quotient: ") (display q) (newline)
    (display "Remainder: ") (display r) (newline)))

;; Using let-values (derived form)
(let-values (((q r) (quotient-remainder 17 5)))
  (list q r))
\end{lambdustcode}
\end{example}

\section{Compatibility Guarantees}
\label{sec:compatibility}

\begin{compatibility}[R7RS Semantic Compatibility]
All R7RS-small expressions have identical semantics in Lambdust when run in Dynamic mode. This includes:
\begin{itemize}
\item Lexical scoping and variable binding
\item Procedure call evaluation order  
\item Proper tail recursion requirements
\item Continuation semantics
\item Multiple value protocols
\item Numeric tower operations
\end{itemize}
\end{compatibility}

The formal semantics in \chapref{ch:formal-semantics} provide mathematical precision for these compatibility guarantees.